\cotefile{157-5}
\scanpage{2}
\uline{$\overline{\mathrm{XVII}}$} Catégories de Modèles (3) :

Cat. à \unsure{cofibrations} \unsure{etc}
% archive sheet number: 1

\scanpage{3}
% printed form
% Line-printer text, read with the sheet rotated through 180 degrees.
{\tiny
\begin{verbatim}
//*   GUTS      0037       I

 GGGGGGGGGG   UU        UU  TTTTTTTTTTTT   SSSSSSSSSS
GGGGGGGGGGGG  UU        UU  TTTTTTTTTTTT  SSSSSSSSSSSS
GG        GG  UU        UU       TT       SS        SS
GG            UU        UU       TT       SS
GG            UU        UU       TT       SSS
GG            UU        UU       TT        SSSSSSSSS
GG     GGGGG  UU        UU       TT         SSSSSSSSS
GG     GGGGG  UU        UU       TT                SSS
GG        GG  UU        UU       TT                 SS
GG        GG  UU        UU       TT       SS        SS
GGGGGGGGGGGG  UUUUUUUUUUUU       TT       SSSSSSSSSSSS
 GGGGGGGGGG    UUUUUUUUUU        TT        SSSSSSSSSS

****P  START  STC   37  GUTS                            ROOM       10.38.39 AM 06 NOV 82  R6.PR1    SYS JES2  STC   37  START  P****
****P  START  STC   37  GUTS                            ROOM       10.38.39 AM 06 NOV 82  R6.PR1    SYS JES2  STC   37  START  P****
****P  START  STC   37  GUTS                            ROOM       10.38.39 AM 06 NOV 82  R6.PR1    SYS JES2  STC   37  START  P****
****P  START  STC   37  GUTS                            ROOM       10.38.39 AM 06 NOV 82  R6.PR1    SYS JES2  STC   37  START  P****
****P  START  STC   37  GUTS            *               ROOM       10.38.39 AM 06 NOV 82  R6.PR1    SYS JES2  STC   37  START  P****
****P  START  STC   37  GUTS                            ROOM       10.38.39 AM 06 NOV 82  R6.PR1    SYS JES2  STC   37  START  P****

- DERNIERE MISE A JOUR DE JOURNAL : 03-11-82
- ATTENTION -------->
            A COMPTER DU LUNDI 8 NOVEMBRE 1982 A 8 HEURES, DE
            NOUVELLES PROCEDURES DE CONNECTIONS VIA TRANSPAC
            POUR UTS ET TSO SERONT MISES EN PLACE. D'ICI LA,
            CONSULTEZ LE JOURNAL OU FAITES '/TRANSPAC' SOUS GUTS,
\end{verbatim}
}
\handwritten{\S~1. \unsure{Andouilleries} préliminaires \hfill 1--24.}

\handwritten{\S~2 $C_I\subset\unsure{\mathrm{Fl}|\mathrm{Fib}}(M(I^\circ))$, pour $I$ ens. \unsure{ordonné} \hfill 25--38
(\unsure{th. principal}, p.~25)}

\handwritten{\S~3. Cat de modèles à cofibrations \hfill 39--\deleted{\unsure{9}}\unsure{70}
(\unsure{th. principaux} p.~50, 60). \unsure{Extensions|Extension}
sur $M(I)$, $I$ ens. \unsure{ordonné}}

\handwritten{\S~4 Catégories à cofibrations spéciales, \hfill 71--\deleted{\unsure{1}}\unsure{83}
\unsure{ext.} sur $M(I)$, $I\in\mathrm{Ob}\,\mathrm{Cat}$.}

\handwritten{\S~5 Catégories de Quillen spéciales \hfill 84--97}

\handwritten{\S~6 Construction \unsure{des|de} dérivateurs \hfill 98--
\unsure{via|wo} catégories de Quillen spéciales. (\unsure{Th. principal} p.~98--101)}

% slanted
\margin{\handwritten{\unsure{$y$}}}
% archive sheet number: 2

\scanpage{5}
\mspage{1}
\deleted{\illegible{1}} $M$ stable par $\underline{\mathrm{Hom}}$

\unsure{Triple|Triples} de Quillen $(W,C,F)$

Dans \deleted{\unsure{$\underline{\mathrm{Hom}}(X,Y)$}} $\mathrm{Ho}_W$ il y a produits

Y-a-t-il $\underline{\mathrm{Hom}}_{\deleted{\unsure{W}}}(X,Y)$ \unsure{pour les mêmes} ?

Supposons que $\underline{\mathrm{Hom}}_{\deleted{W}}$\unsure{s internes} existent dans $M$.

$\underline{\mathrm{Hom}}(X,Y)$ \unsure{bifoncteur} en $X,Y$, mais sous quelles conditions compatible \unsure{à un|avec} $W$ ?

Si \deleted{$X$ \unsure{cofib}}, $Y$ fibrant, $\underline{\mathrm{Hom}}(X,Y)$ ``\unsure{le bon}'' ?

$\underline{\mathrm{Hom}}(X,Y)$ \uline{fibrant sur} $e$ ?

\diagram{\begin{tikzcd}[ampersand replacement=\&]
 S \arrow[r,hook,"TC"] \& T
\end{tikzcd}}

\[
\begin{aligned}
\mathrm{Hom}(T,\underline{\mathrm{Hom}}(X,Y)) &\unsure{-|\to} \mathrm{Hom}(S,\underline{\mathrm{Hom}}(X,Y))\\
\mathrm{Hom}(T\times X,Y)&-\mathrm{Hom}(S\times \unsure{X},Y)
\end{aligned}
\]

\margin{\unsure{épi} ?}

Il faudrait que $TC$ stable par produits cartésiens. par $X^*$ (\unsure{cela|elle} \unsure{implique} idem pour $W$ \unsure{lui-\mbox{$=$}}) : \uline{Ex} $TC=\mathrm{Mono}\cap W$, et $W$ stable par produit cartésien.
% slanted
\margin{$*$ \unsure{évt.} pour $X$ fibrant \unsure{seulement} ?}

\[
\begin{aligned}
\mathrm{Hom}_W(\underset{\text{\scriptsize\inserted{si $S$ cofibrant}}}{S},\underline{\mathrm{Hom}}(X,Y))
 &\simeq \text{Classes d'homotopie à g.}\\
 &\quad\text{dans }\mathrm{Hom}(S,\underline{\mathrm{Hom}}(X,Y))\\
 &\quad\simeq\mathrm{Hom}(S\times X,Y).
\end{aligned}
\]
% archive sheet number: 4

\scanpage{6}
\mspage{P 52}
cette ``promenade'' (dans l'étape ``Points de vue et vision'', n$^\circ$ 6), ces théorèmes ne prennent pour moi tout leur sens que par le contexte nourricier d'un grand thème, initié par une de ces ``idées fécondes''. Leur démonstration dès lors découle, comme de source et sans effort, de la nature même, de la ``profondeur'' du thème qui les porte - comme les vagues du fleuve semblent naître en douceur de la profondeur même de ses eaux, sans rupture et sans effort. Je m'exprime dans un sens tout analogue, mais avec d'autres images, dans la note déjà citée ``La mer qui monte...''.

\noindent\rule{2.1cm}{0.4pt}

18.\quad Quand cet ``avant-propos'' a commencé à tourner à la promenade à travers mon oeuvre de mathématicien, avec mon petit topo sur les ``héritiers'' (bon teint) et sur les ``bâtisseurs'' (incorrigibles), a commencé aussi à apparaître un n o m pour cet avant-propos manqué : ce serait ``L'enfant et le bâtisseur''. Au cours des jours suivants, il devenait de plus en plus clair que ``l'enfant'' et ``le bâtisseur'' étaient un seul et même personnage. Ce nom est donc devenu, plus simplement, ``L'enfant bâtisseur''. Un nom, ma foi, qui ne manquait pas d'allure, et tout fait pour me plaire !

Mais voilà que la réflexion fait apparaître que cet altier ``bâtisseur'', ou (plus modestement) l'enfant-qui-joue-à-faire-des-maisons, ce n'était qu'un des visages du fameux enfant-qui-joue, lequel en avait d e u x. Il y a aussi l'enfant-qui-aime-à-explorer-les-choses, à aller fouiner et s'enfouir dans les sables ou dans les vases boueuses et sans nom, les endroits les plus impossibles et les plus saugrenus... Pour donner le change sans doute (ne serait-ce qu'à moi-même...), j'ai commencé par l'introduire sous le nom flamboyant de ``pionnier'', suivi de celui, plus terre-à-terre mais encore auréolé de prestige, d' ``explorateur''. C'était à se demander, entre le ``bâtisseur'' et le ``pionnier-explorateur'', lequel était le plus mâle, le plus alléchant des deux ! Pile ou face ?

Et puis, en y regardant d'un peu plus près, voilà notre intrépide ``pionnier'' qui se trouve finalement être une f i l l e (qu'il m'avait plu d'habiller en garçon) - une soeur des mares, de la pluie, des brumes et de la nuit, silencieuse et quasiment invisible à force de s'effacer dans l'ombre - celle que toujours on oublie (quand on ne fait mine de se gausser d'elle...). Et j'ai bien trouvé moyen moi aussi, pendant des jours et des jours, de l'oublier - de l'oublier doublement, pourrais-je dire : je n'avais voulu voir d'abord que le garçon (celui qui joue à faire des maisons...) - et même quand je n'ai pu m'empêcher, à force, de voir quand même l ' a u t r e , je l'ai vue encore en garçon, elle aussi...

Pour ce qui est du beau nom pour ma promenade, du coup il ne tient plus du tout. C'est un nom tout-en-yang, tout ``macho'', un nom-qui-boite. Pour le faire tenir pas de guingois, il faudrait y faire figurer l ' a u t r e également. Mais, chose étrange, `` l ' a u t r e ''\quad n ' a \quad p a s \quad v r a i m e n t \quad d e
% archive sheet number: 5

\scanpage{7}
\mspage{2}
\diagram{\begin{tikzcd}[ampersand replacement=\&]
 S \arrow[r,arrows={Hooks[right]-}]\arrow[d] \& T\arrow[dl]\\
 U \& {}
\end{tikzcd}}

\diagram{\begin{tikzcd}[ampersand replacement=\&]
 \& \deleted{\unsure{S\times I}} \& {}\\
 S\vee S\arrow[r,hook,"C"] \& I_S\arrow[r,"W"] \& S
\end{tikzcd}}

\diagram{\begin{tikzcd}[ampersand replacement=\&]
 S\times X\vee S\times X\arrow[r,arrows={Hooks[right]-},"C"]\arrow[ddr] \&
 I_S\times X\arrow[r,"W"]\arrow[d,equal] \& S\times X\\
 \& I_{S\times X}\arrow[d] \& {}\\
 \& {} \& {}
\end{tikzcd}}

\margin{Sommes finies \unsure{distributives} p.r. à produits

$C$ stable par produit cart. par $X$ \unsure{fibrants|fibrant}}

On trouve bien
\[
\mathrm{Hom}_W(S,\underline{\mathrm{Hom}}(X,Y))\simeq\mathrm{Hom}_W(S\times X,Y)
\]
où $S\times X$ est le \uline{produit dans $M$}.

C'est trifonctoriel en $S\in M_c$, $Y\in M_f$, $X\in M_f$.

D'autre part, si $X\unsure{\in} M_f$, alors $S\times X$ est \inserted{aussi} le produit dans $\mathrm{Ho}\,M$.

Ainsi, on a

\uline{Proposition} $(\deleted{\unsure{M}}W,C,F)$ triple de $Q$ dans $M$, \unsure{tel que|satisfaisant}

a) $\underline{\mathrm{Hom}}(X,Y)$ existe dans $M_f$ \deleted{\unsure{pour tout} $X,Y$} \inserted{si $X,Y\in M_f$}, et \inserted{somme $S\amalg T$ \unsure{distrib.} p.r. à produit par $X$}

b) Si $X$ dans $M_f$, $S\longmapsto S\times X$ transforme $TC$ en $TC$, $C$ en $C$

Alors $\underline{\mathrm{Hom}}$ \unsure{internes|interne} existe dans $\mathrm{Ho}_W$, et si $X,Y$ sont dans $M_f$, on a $\underline{\mathrm{Hom}}_{\mathrm{Ho}_W}(X,Y)=\unsure{\mathrm{hot}|\mathrm{Hot}}_W(\underline{\mathrm{Hom}}_M(X,Y))$.
% archive sheet number: 6

\scanpage{8}
\mspage{P 51}
intense. L'un et l'autre sont présents en tout moment du travail, et il y a ce mouvement constant de ``va-et\textsuperscript{\unsure{''}}-vient'' entre les ``temps'' où l'un prédomine, et ceux où prédomine l'autre.

\noindent\rule{2.1cm}{0.4pt}

l'autre mode est présent. Quand je construis, aménage, ou que je déblaie, nettoie, ordonne, c'est le ``mode'' ou le ``versant'' ``yang'', ou ``masculin'' du travail qui donne le ton. Quand j'explore à tâtons l'insaisissable, l'informe, ce qui est sans nom, je suis le versant ``yin'', ou ``féminin'' de mon être.

Il n'est pas question pour moi de vouloir minimiser ou renier l'un ou l'autre versant de ma nature, essentiels l'un et l'autre - le ``masculin'' qui construit et qui engendre, et le ``féminin'' qui conçoit, et qui abrite les lentes et obscures gestations. Je ``suis'' l'un et l'autre - ``yang'' et ``yin'', ``homme'' et ``femme''. Mais je sais aussi que l'essence la plus délicate, la plus déliée dans les processus créateurs se trouve du côté du versant ``yin'', ``féminin'' - le versant humble, obscur, et souvent de piètre apparence.

C'est ce versant-là du travail qui, depuis toujours je crois, a exercé sur moi la fascination la plus puissante. Les consensus en vigueur m'encourageaient pourtant à investir le plus clair de mon énergie dans l'autre versant, dans celui qui s'incarne et s'affirme dans des ``produits'' tangibles, pour ne pas dire finis et achevés - des produits aux contours bien tranchés, attestant de leur réalité avec l'évidence de la pierre taillée...

Je vois \qquad bien, avec le recul, comment ces consensus ont pesé sur moi, et aussi comment j'ai ``accusé le poids'' - en souplesse ! La partie ``conception'' ou ``exploration'' de mon travail était maintenue à la portion congrue jusqu'au moment encore de mon départ, soit. Et pourtant, dans ce coup d'oeil rétrospectif sur ce que fut mon oeuvre de mathématicien, il ressort avec une évidence saisissante que ce qui fait l'essence et la puissance de cette oeuvre, c'est bien ce versant de nos jours négligé, quand il n'est objet de dérision ou d'un condescendant dédain : celui des `` i d é e s '' , voire celui du `` r ê v e '' , nullement celui des ``résultats''. Essayant dans ces pages de cerner ce que j'ai apporté de plus essentiel à la mathématique de mon temps, par un regard qui embrasse une forêt, plutôt que de s'attarder sur des arbres - j'ai vu, non un palmarès de ``grands théorèmes'', mais un vivant éventail d'idées fécondes (*), venant concourir toutes à une même et vaste vision.

\noindent\rule{1.8cm}{0.4pt}

(*) Ce n'est pas que ce qu'on peut appeler les ``grands théorèmes'' manquent dans mon oeuvre, y compris des théorèmes qui résolvent des questions posées par d'autres que moi, que personne avant moi n'avait su résoudre. (J'en passe en revue certains dans la note de b. de p. (***) page 554, de la note ``La mer qui monte...'' ( ReS III, n$^\circ$ 122).) Mais, comme je l'ai souligné déjà dès les débuts de
% archive sheet number: 7

\scanpage{9}
\mspage{3}
\deleted{\uline{Cor.} Si $\unsure{\mathrm{Hom}|\underline{\mathrm{Hom}}}(X,Y)$ existe dans $M$ dès que $Y\in M_f$,}

\deleted{et si}

\uline{Cor.} Plus gén, \deleted{\unsure{pour $X$}} sous l'hyp. \deleted{\unsure{outre $M$, tel}} que la distributivité des sommes p.r. au produit
$\underline{\mathrm{Hom}}_W(X,Y)$ existe dès que $X$ satisfait aux conditions suivantes :
$S\longmapsto S\times X$ transforme $TC$ en $TC$, $C$ en $C$ et \deleted{\unsure{$Y$} il existe $Y\longrightarrow Y'$}
$Y$ est isom. dans $\mathrm{Ho}_W$ à un $Y'\in M_f$ tel que
$\underline{\mathrm{Hom}}_M(X,Y')$ existe dans $M$. \unsure{On a}, pour un tel couple $(X,Y)$,
$\underline{\mathrm{Hom}}_W(X,Y)\simeq\unsure{\mathrm{hot}|\mathrm{Hot}}_W(\underline{\mathrm{Hom}}_M(X,Y'))$

\uline{Corollaire} 1 Existence des $\underline{\mathrm{Hom}}_W$ dans $\mathrm{Hot}_W$, quand $W$ satisfait $\mathrm{Loc}(1\mathrm{bis},2,3,4,7\mathrm{bis})$

$\mathrm{Hom}$ internes dans $\mathrm{Hot}_W(S)$ ? O.K.,

car \deleted{si $C=\mathrm{Mono}$,}

% slanted
\margin{\uline{NB} \unsure{on|--} \unsure{généralise|généraliser} à un produit $\otimes$ au lieu du produit cartésien... ?}

\uline{Corr. 2} Supposons que l'on ait $\unsure{M|CM}$ satisfaisant \inserted{$+$ distr. des sommes p.r. aux produits} a) b) c) d) e) \deleted{f) g)} (pr. \unsure{de mon topo}), et $W$ satisfaisant
\inserted{a) $W\cap\mathrm{Mono}$ stable par \unsure{coch.} de base, b) $\unsure{(\mathrm{Mono})_{\mathrm{fib}}}\subset W$}
\deleted{$(W,C,F)$ \unsure{couple de Quillen clos}, avec $C=\mathrm{Mono}$,}
c) $W$ commute aux $\varinjlim$ \unsure{cofiltrantes|filtrantes}, et accessible. Supposons de plus que $M$ soit stable par $\underline{\mathrm{Hom}}$ internes. \deleted{\unsure{A}} \deleted{que les} \uline{Ces conditions sont stables} par passage de $(M,W)$ à $(M(I),W(I))$, où $I$ dans Cat. Elles impliquent que $(W,\mathrm{Mono},(\mathrm{Mono}\cap W)_{\unsure{*}})$ est un triple de Quillen \inserted{clos} (et idem pour $W(I)$), et que les
% archive sheet number: 8

\scanpage{10}
\mspage{P 50}
profonde à installer une pièce après l'autre, et à voir peu à peu s'instaurer, parmi ces salles, ces chambres et ces réduits l'ordre harmonieux de la maison vivante - belle, accueillante, bonne pour y vivre. Car l a \quad m a i s o n, avant tout et secrètement en chacun de nous, c'est aussi l a \quad m è r e - ce qui nous entoure et nous abrite, à la fois refuge et réconfort ; et peut-être (plus profondément encore, et alors même que nous serions en train de la construire de toutes pièces) c'est cela aussi dont nous sommes nous-mêmes issus, ce qui nous a abrité et nourri, en ces temps à jamais oubliés d'avant notre naissance... C'est aussi l e \quad G i r o n .

Et l'image apparue spontanément tantôt, pour aller au delà de l'appellation prestigieuse de ``pionnier'', et pour cerner la réalité plus cachée qu'elle recouvrait, était elle aussi dépouillée de tout accent ``héroïque''. Là encore, c'était l'image archétype du maternel qui est apparue - celle de la ``matrice'' nourricière et de ses informes et obscurs labeurs...

Ces deux pulsions qui m'apparaissaient comme ``de nature différente'' sont finalement plus proches que je ne l'aurais pensé. L'une et l'autre sont dans la nature d'une `` p u l s i o n \quad d e \quad c o n t a c t '' , nous portant à la rencontre de `` l a \quad M è r e '' : de Celle qui incarne e t ce qui est proche, ``connu'', e t ce qui est ``inconnu''. M'abandonner à l'une ou l'autre pulsion, c'est `` r e t r o u v e r \quad l a \quad M è r e ''. C'est renouveler le contact à la fois au p r o c h e , au ``plus ou moins connu'', et au `` l o i n t a i n '', à ce qui est ``inconnu'' mais en même temps pressenti, sur le point de se faire connaître.

La différence ici est de tonalité, de dosage, non de nature. Quand je ``bâtis des maisons'', c'est le ``connu'' qui domine, et quand ``j'explore'', c'est l'inconnu. Ces deux ``modes'' de découverte, ou pour mieux dire, ces deux aspects d'un même processus ou d'un même travail, sont indissolublement liés. Ils sont essentiels l'un et l'autre, et complémentaires. Dans mon travail mathématique, je discerne un mouvement de va-et-vient constant entre ces deux modes d'approche, ou plutôt, entre les moments (ou les périodes) où l'un prédomine, et ceux où prédomine l'autre (*). Mais il est clair aussi qu'en chaque moment, et l'un et

\noindent\rule{2.1cm}{0.4pt}

(*) Ce que je dis ici sur le travail mathématique est vrai également pour le travail de ``méditation'' (dont il sera question un peu partout dans Récoltes et Semailles). Il n'y a guère de doute pour moi que c'est là une chose qui apparaît dans tout travail de découverte, y compris dans celui de l'artiste (écrivain ou poète, disons). Les deux ``versants'' que je décris ici peuvent être vus également comme étant, l'un celui de l ' e x p r e s s i o n et de ses exigences ``techniques'', l'autre celui de la r é c e p t i o n (de perceptions et d'impressions de toutes sortes), devenant i n s p i r a t i o n par l'effet d'une attention
% archive sheet number: 9

\scanpage{11}
\mspage{4}
$\mathrm{Hom}$ internes existent dans \deleted{les} $\mathrm{Ho}_W$ et les $\mathrm{Ho}_{W(I)}$, et se calculent comme dit \unsure{dans|sur} la prop.

Le seul problème est l'existence des $\underline{\mathrm{Hom}}$ internes dans les $M(I)$, on aura
\[
\begin{aligned}
\underline{\mathrm{Hom}}_{M(I)}(X,Y)
 &=\bigl(i\longmapsto
 \underbrace{\underset{\substack{\downarrow\\\inserted{\substack{\text{\unsure{à val.} dans}\\M}}}}{\underline{\mathrm{Hom}\,\mathrm{ext}}}
 (X\vert I_{/i},Y\vert I_{/i})}_{\displaystyle
 \subset\prod_{\substack{\deleted{\unsure{(u,j)\in}}\\(\unsure{j,u})\in\mathrm{Ob}\,I_{/i}}}\underline{\mathrm{Hom}}_M(X(j),Y(j))}
\end{aligned}
\]

- - - -

\uline{NB} \inserted{$(X,Y)\mapsto$} $\underline{\mathrm{Hom}}(X,Y)$ définit un morphisme de dérivateurs
\[
\underset{\diagup}{\mathbb{D}^{\circ}}\times\underset{\diagup}{\mathbb{D}}\unsure{-|\to}\underset{\diagup}{\mathbb{D}}
\]
Si $f:I'-I$
\[
 f^*(\underline{\mathrm{Hom}}_{\mathbb{D},I}
 (\underset{\substack{\uparrow\\\inserted{\mathbb{D}(I)}}}{X},\underset{\uparrow}{Y}))
 \xrightarrow{\sim}\underline{\mathrm{Hom}}_{\mathbb{D},I'}(f^*X,f^*Y)
\]
\margin{Non !}

\uline{Plutôt} Si $X_{.}\in\mathbb{D}(I)$, $Y_{.}\in\mathbb{D}(J)$, alors on définit
\[
\begin{aligned}
\underline{\mathrm{Hom}}_{\square}(X_{.},Y_{.})&\in\mathbb{D}(I^\circ\times J)=\mathbb{D}_{I^\circ}(J)\\
&=\mathbb{D}_J(I^\circ)
\end{aligned}
\]
Dont il faudrait donner une car. intrinsèque en termes de $\mathbb{D}$. \deleted{\illegible{1}}
% archive sheet number: 10

\scanpage{12}
\mspage{P 49}
celle du pionnier ou de l'explorateur, parti à la découverte et à l'exploration de mondes inconnus et sans nom, m'appelant sans cesse pour que je les connaisse et les nomme. Cette pulsion-là, et l'énergie que j'y investissais (comme à la dérobée, quasiment !), étaient constamment à la portion congrue.

Pourtant, je savais bien au fond que c'était cette énergie-là, dérobée (pour ainsi dire) à celle que je devais à mes ``tâches'', qui était de l'essence la plus rare et la plus déliée - que la ``création'' dans mon travail de mathématicien, c'était avant tout l à qu'elle se plaçait : dans cette attention intense pour appréhender, dans les replis obscurs, informes et moites d'une chaude et inépuisable matrice nourricière, les premières traces de forme et de contours de ce qui n'était pas né encore et qui semblait m'appeler, pour prendre forme et s'incarner et naître... Dans le travail de découverte, cette attention intense, cette sollicitude ardente sont une force essentielle, tout comme la chaleur du soleil pour l'obscure gestation des semences enfouies dans la terre nourricière, et pour leur humble et miraculeuse éclosion à la lumière du jour.

Dans mon travail de mathématicien, je vois à l'oeuvre surtout ces deux forces ou pulsions, également profondes, de nature (me semble\unsure{-{}-|-}t-il) différentes. Pour évoquer l'une et l'autre, j'ai utilisé l'image du b â t i s s e u r , et celle du p i o n n i e r ou de l'explorateur. Mises côte-à-côte, l'une et l'autre me frappent soudain comme vraiment très ``yang'', très ``masculines'', voire ``macho'' ! Elles ont la résonance altière du mythe, ou celle des ``grandes occasions''. Sûrement elles sont inspirées par les vestiges, en moi, de mon ancienne vision ``héroïque'' du travail créateur, la vision super-yang. Telles quelles, elles donnent une vision fortement\unsure{·} teintée, pour ne pas dire figée, ``au garde à vous'', d'une réalité bien plus fluide, plus humble, plus ``simple'' - d'une réalité v i v a n t e .

Dans cette mâle pulsion du ``bâtisseur'', qui semble sans cesse me pousser vers de nouveaux chantiers, je discerne bien pourtant, en même temps, celle du c a s a n i e r : de celui profondément attaché à `` l a '' maison. Avant toute autre chose, c'est `` s a '' maison, celle des `` p r o c h e s '' - le lieu d'une intime entité vivante dont il se sent faire partie. Ensuite seulement, et à mesure que s'élargit le cercle de ce qui est ressenti comme ``proche'', est-elle aussi une ``maison pour tous''. Et dans cette pulsion de ``faire des maisons'' (comme on ``ferait'' l'amour...) il y a aussi et avant tout une t e n d r e s s e. Il y a la pulsion du c o n t a c t avec ces matériaux qu'on façonne un à un, avec un soin amoureux, et qu'on ne connaît vraiment que par ce contact aimant. Et, une fois montés les murs et posés les poutres et le toit, il y a la satisfaction
% archive sheet number: 11

\scanpage{14}
\fbox{\S 1} \uline{Préliminaires}
\margin{$\longleftarrow\quad 1,\unsure{2}\quad !!\longrightarrow$}
% slanted
\margin{\uline{NB.} Les préliminaires se \unsure{compliquent|complètent} \unsure{et se poursuivent} \unsure{dans le paragr. 2 jusqu'en} \illegible{1} \unsure{et|p.} \unsure{$\leftarrow$}}
% slanted
\margin{\deleted{\unsure{Il suffit que} \unsure{il soit défini pour les} limites \unsure{finies}, \unsure{sommes} \unsure{et} \illegible{2} \unsure{deux objets}}}

\uline{Proposition 1} Soit $M$ une catégorie stable \unsure{pour} $\unsure{\varinjlim}$, \unsure{quelconque|quelconques}, $C\subset\mathrm{Fl}\,M$ contenant les iso, stable par \unsure{cochgt.|chgt.} de base et par composition \uline{transfinie}. Soit, pour tout \inserted{petit} ens. ordonné $I$, $C_I\subset\mathrm{Fl}\, M(I^\circ)$ l'ens des flèches \inserted{$u_.:X_.\unsure{-|\to}Y_.$} telles que $\forall i_0\in I$, \deleted{considérant} \inserted{dans} le diagramme

\diagram{\begin{tikzcd}[ampersand replacement=\&,column sep=large,row sep=large]
 \&X'_{i_0}=\varinjlim_{i<i_0}X_i
 \arrow[r]\arrow[d,"{\alpha_{i_0}=\lim_{i<i_0}u_i}"']\arrow[dr,phantom,"\text{cocart}"]
 \& X_{i_0}\arrow[d,"{\unsure{\beta}_{i_0}}"']\arrow[dd,bend left=55,dashed,"{u_{i_0}}"]\\
 (\alpha)\&Y'_{i_0}=\underset{i<i_0}{\underline{\lim}}\,Y_i\arrow[r]
 \&\widetilde{Y}_{i_0}\arrow[d,"{\widetilde{u}_{i_0}}"']\&\text{(défini comme somme amalgamée)}\\
 \&{}\&Y_{i_0}
\end{tikzcd}}
% slanted
\margin{Condition stable par passage de $I$ à un sous-ens. ordonné \uline{ouvert} $I'$ de $I$}

\unsure{la flèche|les flèches} $\widetilde{u}_{i_0}$ \unsure{soit|soient} dans $C$. \deleted{Soit alors $I$ \unsure{soit}} \inserted{Je dis que} \deleted{alors} si $I$ est \uline{bien ordonné}, \deleted{\unsure{et}} alors
$u_.\in C_I$ implique a) que $\varinjlim u_.:\varinjlim X_i-\varinjlim Y_i$ est dans $C$, et b) $\forall i\in\deleted{\unsure{C}} I$, $(u_i:X_i-Y_i)\in C$.

% slanted
\margin{Dire aussi que l'on a b$'$) Si $i'\leq i$, alors \unsure{dans le} diagramme \unsure{similaire|tributaire} à \unsure{$(\alpha)$}
\diagram{\begin{tikzcd}[ampersand replacement=\&,column sep=tiny,row sep=tiny]
 X'_{i'}\arrow[r]\arrow[d] \& X'_i\arrow[d]\\
 Y'_{i'}\arrow[r,-]\arrow[ur,phantom,"\text{cocart}" description] \&\widetilde{Y}'_i\arrow[d,"{u_{i',i}}"]\\
 \&Y'_i
\end{tikzcd}}
on a $u_{i',i}\in C$ \unsure{ce qui} généralise a) \unsure{s'obtient}, retrouve cas en \unsure{prenant} $i'=i_0,\ i=i_0+1)$ le cas b) s'obtient en faisant $i'=0$, $i$ quelconque, et l'hyp. $(X'_{i_0}\to Y'_i)\in C$ en \unsure{supposant} $\widetilde{u}_i:X_i\to Y_i$ \unsure{un} \illegible{1} $i$, \unsure{afin} $\in C$.}

\uline{Cor.} \deleted{\unsure{Démonstration}} Soit $W$ \inserted{$\subset\mathrm{Fl}\,M$, \uline{localiseur}} modérément saturé, tel que $\overset{\inserted{TC\overset{\mathrm{déf}}{=}}}{C\cap W}$ stable par coch. de base, et par composition \uline{transfinie}. Soit $u_.:X_.-Y_.$ dans $C_I$.

Alors $u_.\in W_{I^\circ}$ \deleted{\unsure{(I)}} (i.e. $\forall i\in I$, on a $u_i\in W$) ssi $u_.\in(TC)_I$, i.e. ssi pour \deleted{les} $i\in I$, la flèche $\widetilde{u}_i$ de $(\alpha)$ est dans $TC$, i.e. \deleted{ce qui}
% archive sheet number: 13

\scanpage{15}
\mspage{PU 37}
\deleted{traditionnellement, de paire d' o p p o s é s , comme c'est le cas pour les couples voisins
\[
\text{le précis }\handwritten{\longrightarrow}\text{ le vague}\qquad\text{ou}\qquad
\text{le particulier }\handwritten{\longrightarrow}\text{ le général , -\unsure{.}}
\]
et p\deleted{a}\inserted{\handwritten{\unsure{ou}}}r les quatre autres couples qui suivent dans notre zig-zag. Je présume que même de nos jours, il y aura peu de scientifiques qui songeraient à récuser ce couple, comme exprimant l'harmonie idéale recherchée dans la démarche scientifique.

Pour ma part, si quelqu'un s'était avisé de me poser la question, nul doute que je me serais reconnu dans cet idéal-là, sans y réfléchir à deux fois. Aujourd'hui encore je ne le renierais pas. Mais, comme tout idéal, celui-ci ne touche que la pure surface des choses. Ce n'est pas dans l'idéal, mais dans la riche r é a l i t é et dans sa mère, le r ê v e , que réside la profondeur, et la fécondité qui est sienne.

13.\quad Les épousailles de l'ordre et du mystère

(23 mars)\quad La nuit dernière, en suivant les associations de couples yin-yang suscitées par la réflexion, j'ai vu se prolonger encore sur la gauche notre joli zig-zag de hier, de façon à entrer dans des types de qualités qui ne sont pas particulières seulement à la pensée. Dans le diagramme initial, j'ai remplacé les termes ``le général'' (donc des adjectifs ``substantivisés'') par les noms correspondants ``généralités'', etc. Cela conduit alors à le prolonger comme suit (*) :

\diagram{\begin{tikzcd}[ampersand replacement=\&,column sep={6.7mm,between origins},row sep=10pt,cells={nodes={font=\tiny,inner sep=0.5pt}}]
{} \& \text{discipline}\handwritten{\arrow[dl]\arrow[dr,line width=.52pt]} \& {} \& \text{contrôle}\handwritten{\arrow[dl]\arrow[dr,line width=.52pt]} \& {} \& \text{volonté}\handwritten{\arrow[dl]\arrow[dr,line width=.52pt]} \& {} \& \text{rigueur}\handwritten{\arrow[dl]\arrow[dr,line width=.52pt]} \& {} \& \text{précision}\handwritten{\arrow[dl]\arrow[dr,line width=.52pt]} \& {} \& \text{clarté}\handwritten{\arrow[dl]\arrow[dr,line width=.52pt]} \& {} \& \text{savoir}\handwritten{\arrow[dl]\arrow[dr,line width=.52pt]} \& {}\\
\text{imagination} \& {} \& \text{jeu} \& {} \& \text{abandon} \& {} \& \text{spontané\deleted{\unsure{te}}\handwritten{\unsure{it}}é} \& {} \& \text{généralité} \& {} \& \text{flou} \& {} \& \text{obscurité} \& {} \& \text{mystère ,}
\end{tikzcd}}

où la ligne du haut est encore formée des termes yang, celle du bas des termes yin.

\noindent\rule{2.1cm}{0.4pt}

(*) (25 mars)\quad Avant de prolonger le zig-zag de la veille, j'avais été conduit à remplacer le terme yang ``le particulier'' (s'accouplant avec ``le général'' ou ``généralité'', et faisant presque double-emploi avec son voisin de droite ``précision'') par le terme ``rigueur'', pour former le nouveau couple
\[
 \text{rigueur - généralité}
\]
venant parfaire son voisin (ou ``concubin'') de droite
\[
 \text{précision - généralité .}
\]
C'est bien le premier de ces deux couples qui me paraît le mieux ``assorti''. S'il est vrai que souvent, on perd en précision ce qu'on gagne en généralité et inversement, une telle situation ne se présente j a m a i s pour le couple ``rigueur -}
% archive sheet number: 14

\scanpage{16}
\mspage{3}
revient au même, dans $W$ (puisque par hypothèse elle est déjà dans $C$). \deleted{\unsure{Or}}\unsure{Si} cette condition est satisfaite, alors $\varinjlim u_i:\varinjlim X_i\longrightarrow\varinjlim Y_i$ est dans $W$.

Dém. du corollaire. Supposons $u_.\in(TC)_I$, alors par la proposition, $\varinjlim_{i\in I}u_i$ est dans $TC$, \deleted{\unsure{donc dans $W$}} \inserted{ou ce qui revient} au $=$ dans $W$. D'autre part, appliquant ce $=$ résultat aux syst. inductifs induits sur $I_{<i_0}$, on trouve que dans $(\alpha)$, la \deleted{\unsure{$\beta$}} flèche \deleted{$\unsure{\lim}_{i<i_0}u_i$} $\alpha_{i_0}$ est dans $TC$, donc aussi $\beta_{i_0}$ par cochgt. de base, donc aussi $u_{i_0}=\widetilde{u}_{i_0}\beta_{i_0}$. \deleted{\unsure{Donc}} Donc il reste à prouver qu'inversement, si on suppose les $u_{i_0}\in W$, alors les $\widetilde{u}_{i_0}$ sont eux aussi dans $W$. (i.e. dans $TC=C\cap W$), ou ce qui revient au $=$, que les $\beta_{i_0}$ sont dans $W$, ou encore dans $TC$. Mais comme $TC$ est stable par coch. de base, il suffit \inserted{pour celà} de prouver que $\alpha_{i_0}\in TC$, ou encore $\alpha_{i_0}\in C$. Donc on est ramené à prouver que \deleted{\unsure{si} $u_.\in W_{I^\circ}$}
\[
 u_.\in W_{I^\circ}\cap C_I\Longrightarrow \left(\varinjlim_{\unsure{i}<i_0}u_i\right)\in W \tag*{$c_{i_0})$}
\]
pour \unsure{tt} $i_0\in I$, \unsure{ou}, changeant de notations, à prouver

Nous prouverons ensemble a) b) c) par induction transfinie sur le nb ordinal $\alpha$, tel que $I\simeq I_\alpha=\{\text{ens des ordinaux }i<\alpha\}$.
% archive sheet number: 15

\scanpage{17}
\mspage{4}
\deleted{(Le cas $I=\varnothing$ étant trivial).}

$1^\circ$) Cas $\alpha=0$ i.e. $I$ \deleted{réduit à un élément, \unsure{donc} $M(I^\circ)\overset{\mathrm{can}}{\simeq} M$. Alors \unsure{il est tout} immédiat que $C_I\unsure{\simeq|=}C$, $(TC)_I=TC$, $W_{I^\circ}=W$, et l'énoncé est trivial, tautologique.} \inserted{$=\varnothing$ : il est trivial}

$2^\circ$) Si c'est OK pour $I_\alpha$, aussi pour $I_{\alpha+1}$ \inserted{$=\{\beta\mid\beta\leq\alpha\}$}, \deleted{\unsure{Ici c'est le plus délicat, d'abord}} \inserted{$I_{\alpha+1}$ ayant un plus grand élément $\alpha$,} a) est contenu dans b) (pour $i=\alpha$, \unsure{à|=} plus \deleted{\unsure{grand}} élément de $I_{\alpha+1}$).

\deleted{Considérons alors le diagramme} \deleted{Alors c) résulte de l'hypothèse de récurrence, qui implique}
Pour prouver b), il reste à prouver que $u_\alpha\in C$. On regarde le diagramme $(\alpha)$ \inserted{pour $i_0=\alpha$.} Par \deleted{\unsure{ailleurs}} l'hyp. de récurrence \inserted{(pour a)} appliquée à $I_\alpha$, on a $\alpha_{\unsure{i_0|\alpha_0}}\in C$, donc par cochgt. de base $\beta_{i_0}\in C$, donc par composition $u_{i_0}=\widetilde{u}_{i_0}\beta_{i_0}\in C$.

Pour prouver c), il reste à prouver que $\varinjlim_{i<\alpha}u_i\in W$. \deleted{mais par hyp. de réc., \unsure{on} sait que les $u_i$ $(i<\alpha)$ sont dans $W$, \unsure{donc}} Mais par l'hyp. de réc. appliquée à $I_\alpha=\{i\mid i<\alpha\}$, on sait déjà que $(u_i)_{i<\alpha}\in (TC)_{I_\alpha}$, donc par la partie prouvée du corollaire ou de la prop., $\varinjlim_{i<\alpha}u_i\in TC$.

$3^\circ$) Si $\alpha$ un ordinal limite, et si a) b) c) vrais pour $I_\beta$ \unsure{avec} $\beta<\alpha$, aussi pour $I_\alpha$.

- Cette fois, b) est contenu dans l'hyp. de
% archive sheet number: 16

\scanpage{18}
\mspage{5}
récurrence, il faut prouver a) i.e. \deleted{\unsure{l'iso}} que
\deleted{$\underline{\lim} X_i\to\unsure{\lim}$}
\[
\begin{aligned}
u_\alpha\overset{\mathrm{déf}}{=}\varinjlim_{i<\alpha}\deleted{\unsure{X}}u_i:\quad&
 \underset{\inserted{\substack{\mathrm{déf}\ \Vert\\\varinjlim_{i<\alpha}X_i}}}{X_\alpha}
 \unsure{-|\to}
 \underset{\inserted{\substack{\Vert\\\varinjlim_{i<\alpha}Y_i}}}{Y_\alpha}
\end{aligned}
\]
est dans $C$. \deleted{Et} Quant \inserted{à c, il est déjà} \deleted{\unsure{si on suppose, en} plus de $u_.\in C_{I_\alpha}$, que $u_.\in W_{I^\circ}$, donc (par hyp. de réc.) que pour tt $i_0<\alpha$, dans $(\alpha)$ on a
\[
 \widetilde{u}_{i_0}\in TC\overset{\mathrm{déf}}{=}C\cap W
\]
(et non seulement $\widetilde{u}_{i_0}\in C$), alors on a $u_\alpha\in$}
contenu dans l'hyp. de récurrence.

Donc il reste finalement à prouver
\[
 u_\alpha\in C
\]
quand on suppose que la prop. est prouvée pour les ordinaux $\beta<\alpha$.

On va construire un troisième système \unsure{ordinal|ordonné} sur $I_\alpha$, en \deleted{\unsure{posant}} introduisant pour tt $0\leq i<\alpha$, \inserted{i.e. $i\in I_\alpha$,} $Z_i$ par le carré cocartésien
\diagram{\begin{tikzcd}[ampersand replacement=\&]
 X_i\arrow[r,-]\arrow[d]\arrow[dr,phantom,"\text{cocart}" description] \& X_\alpha\arrow[d]\\
 Y_i\arrow[r] \& Z_i
\end{tikzcd}}
les flèches de transition entre les $Z_i$ sont évidentes, et on aura \deleted{\unsure{c'est}} $\varinjlim_i Z_{\deleted{\unsure{\alpha}}i}\simeq Y_\alpha$. Ainsi $u_\alpha:X_\alpha\inserted{(=Z_{-1})}-Y_\alpha$ \deleted{\unsure{se}} est devenu un composé transfini de morphismes \deleted{\unsure{$Z_i\to Z_{i+1}$} \unsure{$(i\leq\alpha)$} \illegible{1}}
% archive sheet number: 17

\scanpage{19}
\mspage{6}
\[
 Z_i\overset{v_{i+1}}{\unsure{-|\longrightarrow}}Z_{i+1}\qquad(-1\leq i<\alpha) \tag*{$(*)$}
\]
et
\[
 \varinjlim_{i<\deleted{\unsure{\beta}}i_0}Z_i\xrightarrow{v_{i_0}}Z_{\deleted{\beta}i_0}\deleted{\unsure{\text{est}}}\qquad
 \deleted{\beta}i_0\text{ ordinal limite }<\alpha \tag*{$(**)$}
\]
\deleted{\unsure{cette} \unsure{fois}}

(NB le cas où \deleted{$\beta$}\inserted{$i_0$} n'est \uline{pas} ordinal limite n'est autre que \unsure{$(*)$} pour $i=i_0-1$) Tout revient à prouver que ces flèches sont $\in C$. Alors leur composé transfini \inserted{aussi}. \deleted{Cas $(**)$, où $\unsure{\beta|i}_0\leq\alpha$ \unsure{qui jamais limite}}

Cas $(*)$
% slanted
\margin{Pour se donner confiance -- on pourrait traiter \deleted{\unsure{de}} de suite le cas général}
\diagram{\begin{tikzcd}[ampersand replacement=\&]
 X_i\arrow[r]\arrow[d,"u_i"']\arrow[dr,phantom,"\text{cocart}" description] \& X_{i+1}\arrow[r]\arrow[d,"\beta_{i+1}"']\arrow[dr,phantom,"\text{cocart}" description] \& X_\alpha\arrow[d]\\
 Y_i\arrow[r]\&\widetilde Y_{i+1}\arrow[r]\arrow[d,"\widetilde u_{i+1}"']\arrow[dr,phantom,"\text{cocart}" description] \& Z_i\arrow[d,"v_i"]\\
 \&Y_{i+1}\arrow[r]\& Z_{i+1}
\end{tikzcd}}
Comme $\widetilde u_{i+1}\in C$, $v_i\in C$ par \unsure{cochangement|changement} de base, OK.

Il faut cependant traiter à part le cas \uline{$i=-1$}, où la flèche $X_\alpha\to Z_0$ que cela induit
\diagram{\begin{tikzcd}[ampersand replacement=\&]
 X_0\arrow[r]\arrow[d,"u_0"']\arrow[dr,phantom,"\text{cocart}" description] \& X_\alpha=Z_{-1}\arrow[d,"v_{-1}"]\\
 Y_0\arrow[r]\& Z_0
\end{tikzcd}}
où $u_0\in C$, donc $v_{-1}\in C$, OK.

Cas $(**)$. Si \deleted{\unsure{$\beta$ est}} $i_0<\alpha$ est ordinal limite, \unsure{transportons|inspectons} \unsure{carr. bas|cas, bas} -- par \unsure{passage} à la limite
\[
 Z'_{i_0}\overset{\mathrm{déf}}{=}\varinjlim_{i<i_0}Z_i\simeq
 \deleted{\unsure{X}}Y'_{\deleted{\beta}i_0}\amalg_{X'_{i_0}}X_\alpha ,
\]
\inserted{(cf diagramme $(\alpha)$ pour déf. de $X'_{i_0}$, $Y'_{i_0}$)} on a le diagramme suivant \uline{greffé} \unsure{sous|sur} $(\alpha)$
\deleted{\diagram{\begin{tikzcd}[ampersand replacement=\&]
 X'_{i_0}\arrow[r,-]\arrow[d,-]\arrow[dr,phantom,"\text{cocart}" description]\& X_{i_0}\arrow[r]\arrow[d]\& X_\alpha\\
 Y'_{i_0}\arrow[r,-]\&\widetilde Y_{i_0}\arrow[d]\&\\
 \&{}\&
\end{tikzcd}}}
% archive sheet number: 18

